\documentclass[aspectratio=169]{beamer}
\usetheme{Madrid}
\usecolortheme{default}

% Packages
\usepackage{graphicx} % For images
\usepackage{booktabs} % For professional tables
\usepackage{amsmath, amssymb} % For math equations
\usepackage{microtype} % Typography polishing

% Custom Palette & Styling
\definecolor{AmritaBlue}{RGB}{0, 51, 102}
\definecolor{AccentGreen}{RGB}{34, 139, 34}
\setbeamercolor{structure}{fg=AmritaBlue}
\setbeamercolor{block title}{bg=AmritaBlue, fg=white}
\setbeamercolor{block body}{bg=gray!10}
\setbeamercolor{alertblock title}{bg=red!80!black, fg=white}
\setbeamercolor{alertblock body}{bg=red!5}

% Title Page Information
\title[Alzheimer's Detection via Speech Processing]{Alzheimer's Disease Detection Using Speech Processing with Explainable AI}
\subtitle{Mid-Semester Project Review}
\author[Divagar, Abhishek, Prem]{
    \textbf{Divagar S} [CB.SC.U4AIE23223] \\
    \textbf{Abhishek Pandey} [CB.SC.U4AIE23205] \\
    \textbf{Prem} [CB.SC.U4AIE23241]
}
\institute[Amrita Vishwa Vidyapeetham]{22AIE450 Speech Processing \\ Department of AI \& Data Science}
\date{August 2026}

\begin{document}

% ===== Slide 1: Title Slide =====
\begin{frame}
    \titlepage
\end{frame}

% ===== Slide 2: Agenda =====
\begin{frame}{Agenda}
    \tableofcontents
\end{frame}

% ===== Slide 3: Introduction and Motivation =====
\section{Introduction}
\begin{frame}{Introduction and Motivation}
    \begin{block}{Background Context}
        Alzheimer's Disease (AD) is a progressive brain disorder that gradually affects memory, thinking skills, and spoken language.
    \end{block}
    \begin{itemize}
        \item \textbf{Motivation:} Brain scans (MRI/PET) and clinical visits are expensive, time-consuming, and hard to access. Analyzing spontaneous speech provides a non-invasive, quick, and low-cost alternative for early screening.
        \item \textbf{Speech Indicators:} Subtle changes like longer pauses, unstable pitch, slower speech rate, and reduced voice clarity can indicate early cognitive decline.
        \item \textbf{Need for Explainability:} Adding Explainable AI (XAI) helps doctors understand why a model flagged a patient, building trust in model predictions.
    \end{itemize}
\end{frame}

% ===== Slide 4: Problem Definition =====
\section{Problem Definition}
\begin{frame}{Problem Definition}
    \begin{alertblock}{Main Challenges in AD Detection}
        Current diagnostic methods are costly and subjective, while standard AI models often act as unexplainable black boxes.
    \end{alertblock}
    
    \begin{enumerate}
        \item \textbf{High Cost and Limited Access:} MRI/PET brain scans are expensive and not easily available in primary healthcare clinics.
        \item \textbf{Subjective Cognitive Tests:} Standard written questionnaires (like MMSE) depend on clinician evaluation and take time to conduct.
        \item \textbf{Limitations of Basic Audio Features:} Traditional features like pitch or standard MFCCs alone cannot capture complex speech patterns associated with Alzheimer's.
        \item \textbf{Lack of Model Transparency:} Machine learning models must provide understandable explanations before they can be used in medical applications.
    \end{enumerate}
\end{frame}

% ===== Slide 5: Project Overview =====
\begin{frame}{Project Architecture \& Workflow}
    \begin{center}
        \textbf{\Large System Pipeline Overview}
    \end{center}
    \vfill
    \begin{block}{Pipeline Flow (Mid-Term Stage)}
    \begin{itemize}
        \item \textbf{Speech Input:} Spontaneous audio recordings (.wav format) from the ADReSSo 2021 dataset.
        \item \textbf{Audio Preprocessing:} Converting to 16kHz, noise reduction, and removing silent gaps using voice activity detection.
        \item \textbf{Feature Extraction:} Extracting 123 physical acoustic features (MFCCs, Pitch, HNR) + Whisper deep audio embeddings.
        \item \textbf{Feature Fusion:} Combining acoustic features and deep embeddings using PCA (reduced to 318 features).
        \item \textbf{Model Training \& Evaluation:} Training ML and PyTorch DL models using 5-Fold Stratified Cross-Validation.
    \end{itemize}
    \end{block}
\end{frame}

% ===== Slide 6: Gaps Identified =====
\section{Gaps Identified}
\begin{frame}{Gaps Identified in Current Research}
    \begin{alertblock}{Research Gaps}
        Key issues observed in existing literature that our work addresses:
    \end{alertblock}
    \begin{itemize}
        \item \textbf{Lack of Explanation:} Most speech-based AI models operate as black boxes without explaining feature importance to clinicians.
        \item \textbf{Low Accuracy of Simple Acoustic Features:} Official ADReSSo acoustic baselines (eGeMAPS + SVM) achieve only ~65.1\% accuracy.
        \item \textbf{Need for Feature Combination:} Combining hand-crafted acoustic properties with deep audio embeddings improves model sensitivity.
        \item \textbf{Small Dataset Variance:} Small speech corpora require careful 5-fold cross-validation to prevent data leakage and overfitting.
    \end{itemize}
\end{frame}

% ===== Slide 7: Objectives =====
\section{Objectives}
\begin{frame}{Project Objectives}
    \begin{block}{Main Goal}
        To build a speech-based machine learning pipeline for Alzheimer's Detection that outperforms standard baselines and provides clear explanations.
    \end{block}
    
    \textbf{Mid-Term Milestone Objectives (Completed):}
    \begin{enumerate}
        \item Clean and preprocess speech audio (16kHz resampling, noise reduction, silence removal).
        \item Extract 123 physical acoustic features (MFCCs, Pitch statistics, HNR) and Whisper audio features.
        \item Combine features using PCA dimensionality reduction (318 components).
        \item Evaluate Logistic Regression, SVM, Random Forest, PyTorch DNN, and 1D-CNN models using 5-Fold Stratified Cross-Validation.
    \end{enumerate}
\end{frame}

% ===== Slide 8: Methodology =====
\section{Methodology}
\begin{frame}{Detailed Methodology \& Technical Workflow}
    \begin{columns}[T]
        \column{0.52\textwidth}
        \textbf{5-Stage Pipeline Workflow:}
        \begin{enumerate}
            \item \textbf{Audio Preprocessing:} Converting to 16kHz, spectral noise reduction (\texttt{noisereduce}), silence trimming (\texttt{librosa.effects.split}, top\_db=25).
            \item \textbf{Feature Extraction:} 
                \begin{itemize}
                    \item Acoustic (123 features): MFCCs (1-40 + Deltas), Pitch F0 (Mean, Std), HNR.
                    \item Deep Features: Whisper pretrained encoder embeddings.
                \end{itemize}
            \item \textbf{Feature Fusion:} Concatenation $\rightarrow$ StandardScaler $\rightarrow$ PCA (318 dimensions).
            \item \textbf{Model Training:} SVM (RBF), Logistic Regression, Random Forest, PyTorch 3-Layer DNN, 1D-CNN.
            \item \textbf{Evaluation:} 5-Fold Stratified Cross-Validation.
        \end{enumerate}
        
        \column{0.45\textwidth}
        \textbf{System Block Diagram:}
        \begin{center}
            \fbox{\parbox{0.95\textwidth}{\centering
                \small
                \textbf{Speech Recordings (.wav)} \\
                $\downarrow$ \\
                \textbf{16kHz Preprocessing \& VAD} \\
                $\downarrow$ \\
                \textbf{Acoustic (123d) + Whisper (1024d)} \\
                $\downarrow$ \\
                \textbf{StandardScaler + PCA (318d)} \\
                $\downarrow$ \\
                \textbf{ML \& PyTorch DL Models} \\
                $\downarrow$ \\
                \textbf{5-Fold Stratified CV Output}
            }}
        \end{center}
    \end{columns}
\end{frame}

% ===== Slide 9: Dataset Description =====
\section{Dataset}
\begin{frame}{Dataset Description (IS2021 ADReSSo Corpus)}
    \begin{table}
    \centering
    \begin{tabular}{lcccc}
        \toprule
        \textbf{Sub-Dataset} & \textbf{Samples} & \textbf{Class Balance} & \textbf{Audio Specs} & \textbf{Target Task} \\
        \midrule
        Diagnosis Train & 166 Files & 87 AD / 79 Healthy & 16kHz Mono .wav & AD vs Healthy \\
        Progression Train & 73 Files & 15 Decline / 58 Stable & 16kHz Mono .wav & Decline Tracking \\
        Challenge Test & 32 Files & Unlabelled & 16kHz Mono .wav & Final Evaluation \\
        \bottomrule
    \end{tabular}
    \caption{ADReSSo speech dataset details used in our project.}
    \end{table}
    \vfill
    \begin{itemize}
        \item \textbf{Speech Task:} Spontaneous speech recorded during the Cookie Theft picture description test.
        \item \textbf{Data Care:} Original dataset audio files remain untouched. Preprocessed clean files are saved separately in \texttt{processed\_audio/}.
    \end{itemize}
\end{frame}

% ===== Slide 10: Preliminary Results (Part 1) =====
\section{Results}
\begin{frame}{Experimental Results (40--50\% Completion Stage)}
    \begin{block}{5-Fold Stratified Cross-Validation Results}
        \textit{Evaluated on 166 ADReSSo speech samples using PCA feature fusion (318 features)}
    \end{block}
    
    \begin{table}
    \centering
    \small
    \begin{tabular}{lccccc}
        \toprule
        \textbf{Model Architecture} & \textbf{Accuracy (\%)} & \textbf{Precision} & \textbf{Recall} & \textbf{F1-Score} & \textbf{ROC-AUC} \\
        \midrule
        \textbf{Logistic Regression} & \textbf{70.52\% $\pm$ 8.06\%} & \textbf{0.7217} & \textbf{0.7111} & \textbf{0.7117} & \textbf{0.7815} \\
        \textbf{PyTorch DNN (3-Layer)} & \textbf{70.48\% $\pm$ 9.07\%} & \textbf{0.7208} & \textbf{0.7235} & \textbf{0.7209} & \textbf{0.7568} \\
        \textbf{SVM (RBF Kernel)} & \textbf{66.29\% $\pm$ 9.54\%} & 0.6665 & 0.7346 & 0.6950 & 0.7440 \\
        Random Forest & 54.80\% $\pm$ 8.67\% & 0.5566 & 0.6562 & 0.5967 & 0.5285 \\
        PyTorch 1D-CNN & 54.21\% $\pm$ 7.80\% & 0.5644 & 0.6458 & 0.5924 & 0.5732 \\
        \bottomrule
    \end{tabular}
    \end{table}
    
    \begin{block}{Baseline Comparison}
        Our Logistic Regression (\textbf{70.52\%}) and PyTorch DNN (\textbf{70.48\%}) outperform the official ADReSSo acoustic baseline (\textbf{65.1\% accuracy}).
    \end{block}
\end{frame}

% ===== Slide 11: Preliminary Results (Part 2) =====
\begin{frame}{Performance Analysis \& Key Findings}
    \begin{columns}[T]
        \column{0.52\textwidth}
        \textbf{Key Observations from Experiments:}
        \begin{itemize}
            \item \textbf{Strong Separation:} Logistic Regression achieved the best ROC-AUC (\textbf{0.7815}) and F1-score (0.7117), showing solid baseline performance.
            \item \textbf{Neural Network Results:} PyTorch 3-Layer DNN achieved \textbf{70.48\% accuracy} with balanced precision (0.7208) and recall (0.7235).
            \item \textbf{Improvement over Baseline:} Combining acoustic features with Whisper embeddings improved accuracy by \textbf{+5.4\%} over standard acoustic features (65.1\%).
        \end{itemize}
        
        \column{0.45\textwidth}
        \textbf{Generated Results Files:}
        \begin{center}
            \fbox{\parbox{0.95\textwidth}{\centering
                \small
                \textbf{Saved Local Artifacts:} \\
                \vspace{0.2cm}
                • \texttt{model\_comparison\_phase1.csv} \\
                • \texttt{roc\_auc\_curves\_phase1.png} \\
                • \texttt{cm\_logistic\_regression\_phase1.png} \\
                • \texttt{cm\_dnn\_phase1.png} \\
                • \texttt{cm\_svm\_rbf\_phase1.png} \\
                \vspace{0.2cm}
                \textit{(Stored in \texttt{results\_40\_50\_percent\_phase/})}
            }}
        \end{center}
    \end{columns}
\end{frame}

% ===== Slide 12: Interim Summary/Conclusion =====
\section{Summary}
\begin{frame}{Interim Progress Summary}
    \begin{enumerate}
        \item \textbf{Pipeline Completed:} Successfully built the audio cleaning, feature extraction (123 acoustic + Whisper), and PCA feature fusion modules.
        \item \textbf{Good Baseline Performance:} Achieved \textbf{70.52\% accuracy} and \textbf{0.7815 ROC-AUC}, beating the official acoustic baseline (65.1\%).
        \item \textbf{Code Organization:} Structured project into clean python scripts in \texttt{src/} with automated pipeline runner \texttt{run_pipeline.py}, saved locally and on GitHub.
    \end{enumerate}
    \vfill
    \begin{center}
        \begin{block}{Current Project Status}
            \centering
            \textbf{Mid-Term Checkpoint Completed} \\
            Baseline Data Preprocessing and Classifier Models Implemented \& Verified.
        \end{block}
    \end{center}
\end{frame}

% ===== Slide 13: Future Scope =====
\section{Future Scope}
\begin{frame}{Future Scope}
    \begin{table}
    \centering
    \small
    \begin{tabular}{p{0.22\textwidth}p{0.70\textwidth}}
        \toprule
        \textbf{Next Phases} & \textbf{Planned Tasks} \\
        \midrule
        \textbf{Fusion} & Gated Cross-Attention module (dynamically weighting acoustic vs deep features). \\
        \textbf{Multi-Task} & Multi-Task DNN to predict AD classification and MMSE cognitive test scores. \\
        \textbf{XAI} & SHAP feature ranking and LIME explanations to make predictions clear to clinicians. \\
        \textbf{Generalization} & Cross-corpus testing between Diagnosis and Progression datasets. \\
        \textbf{Optimization} & Hyperparameter tuning using Optuna and feature ablation tests. \\
        \bottomrule
    \end{tabular}
    \end{table}
\end{frame}

% ===== Slide 14: References =====
\begin{frame}{References}
    \begin{thebibliography}{9}
        \bibitem{1} Luz, S., Haider, F., de la Fuente, S., Fromm, D., \& MacWhinney, B. (2021). Alzheimer's Dementia Recognition through Spontaneous Speech: The ADReSSo Challenge. \textit{Interspeech 2021}.
        \bibitem{2} Balagopalan, A., Eyre, B., Rudzicz, F., \& Novikova, J. (2020). To BERT or not to BERT: Comparing speech and language-based approaches for Alzheimer's disease detection. \textit{Interspeech 2020}.
        \bibitem{3} Lundberg, S. M., \& Lee, S. I. (2017). A unified approach to interpreting model predictions. \textit{Advances in Neural Information Processing Systems (NeurIPS)}.
    \end{thebibliography}
\end{frame}

% ===== Slide 15: Thank You =====
\begin{frame}{Thank You}
    \begin{center}
        \Huge{\textcolor{AmritaBlue}{Thank You!}}
        
        \vspace{0.5cm}
        \Large{\textbf{Alzheimer's Disease Detection Using Speech Processing with Explainable AI}}
        
        \vspace{0.5cm}
        \normalsize{Team Members: Divagar S, Abhishek Pandey, Prem} \\
        \small{Department of AI \& Data Science | Amrita Vishwa Vidyapeetham}
        
        \vspace{0.8cm}
        \textit{\Large Questions \& Discussion?}
    \end{center}
\end{frame}

\end{document}
